% !TEX root = THESIS_REPORT.tex
\documentclass[11pt,a4paper]{report}

% --- Packages ---
\usepackage[utf8]{inputenc}
\usepackage[margin=1in]{geometry}
\usepackage{amsmath,amssymb,amsfonts}
\usepackage{graphicx}
\usepackage{hyperref}
\usepackage{booktabs}
\usepackage{xcolor}
\usepackage{listings}
\usepackage{float}
\usepackage{enumitem}

\hypersetup{
    colorlinks=true,
    linkcolor=blue,
    citecolor=blue,
    urlcolor=blue
}

% --- Code Listing Syntax Highlighting ---
\definecolor{codegreen}{rgb}{0,0.6,0}
\definecolor{codegray}{rgb}{0.5,0.5,0.5}
\definecolor{codepurple}{rgb}{0.58,0.2,0.82}
\definecolor{backcolour}{rgb}{0.96,0.96,0.96}

\lstdefinestyle{pddlstyle}{
    backgroundcolor=\color{backcolour},   
    commentstyle=\color{codegreen},
    keywordstyle=\color{blue}\bfseries,
    numberstyle=\tiny\color{codegray},
    stringstyle=\color{codepurple},
    basicstyle=\ttfamily\small,
    breakatwhitespace=false,         
    breaklines=true,                 
    captionpos=b,                    
    keepspaces=true,                 
    numbers=left,                    
    numbersep=5pt,                  
    showspaces=false,                
    showstringspaces=false,
    showtabs=false,                  
    tabsize=2
}
\lstset{style=pddlstyle}

% --- Title & Metadata ---
\title{\textbf{Towards an Autonomous Lifecycle Agent for PDDL Domain Modelling}}
\author{\textbf{Master Thesis Report}\\
Course of Study: Computer Science\\
Institute of Architecture of Application Systems, University of Stuttgart}
\date{\today}

\begin{document}

\maketitle

% --- Abstract ---
\begin{abstract}
Creating formal planning domain models in languages such as the Planning Domain Definition Language (PDDL) is a foundational bottleneck in automated planning. It requires significant domain expertise, precise logical formulation, and exhaustive manual debugging to prevent syntax collapse, ungrounded actions, or unreachable goals. While recent advances in Large Language Models (LLMs) have showcased strong natural language translation capabilities, direct one-shot LLM translation frequently yields syntactically malformed or semantically inconsistent PDDL models.

To overcome these limitations, this thesis presents a novel \textbf{Unified Autonomous Lifecycle Agent for PDDL Domain Modelling} built on a \textbf{LangGraph closed-loop state machine}. The framework operationalizes the complete 4-phase \textbf{Knowledge Engineering for Planning and Scheduling (KEPS)} lifecycle:
\begin{enumerate}
    \item \textbf{Phase 1 (Acquisition \& Conceptualization):} Natural language requirement extraction and relational entity mapping.
    \item \textbf{Phase 2 (Formulation \& Encoding):} Intermediate Pydantic Abstract Syntax Tree (AST) schema enforcement (\texttt{DomainBlueprint}, \texttt{ProblemBlueprint}) paired with a deterministic programmatic Lisp compiler.
    \item \textbf{Phase 3 (Verification \& Validation - V\&V):} Symbolic binary Oracles (VAL plan validator and Fast Downward planner) coupled with automated diagnostic fault triage and targeted LLM model repair loops.
    \item \textbf{Phase 4 (Maintenance \& Evolution):} A historical regression suite that enforces backward compatibility during domain evolution.
\end{enumerate}

Empirical evaluation on the \texttt{PDDLLifeCycleBench} dataset across International Planning Competition (IPC) benchmark tasks demonstrates that our closed-loop AST agent achieves a \textbf{100\% Syntax Validity Rate}, a \textbf{100\% Plan Validity Rate}, and a \textbf{100\% Repair Success Rate}, significantly outperforming baseline \texttt{ZERO\_SHOT} and \texttt{UNSTRUCTURED\_RETRY} pipelines. The system provides a fully automated, verifiable, and human-out-of-the-loop framework for planning domain engineering.
\end{abstract}

\tableofcontents
\newpage

% --- Chapter 1: Introduction ---
\chapter{Introduction}

\section{Overview of the State of the Art}
Automated planning represents one of the foundational paradigms of Artificial Intelligence, focusing on the autonomous synthesis of action sequences—termed \textit{plans}—that transition a discrete state-transition system from a known initial world state $S_0$ into a desired target goal state $S_g$ \cite{ghallab2004}. In classical AI planning, the world state is modeled under strict deterministic, static, discrete, and fully observable assumptions. The environment semantics are encapsulated through formal modeling languages, most notably the Planning Domain Definition Language (PDDL), which has served as the de facto standard across the automated planning research community since its inception for the International Planning Competition (IPC) \cite{mcdermott1998}.

Formally, a classical planning domain and instance tuple is expressed as $\Pi = \langle \mathcal{T}, \mathcal{P}, \mathcal{O}, \mathcal{A}, I, G \rangle$, where:
\begin{itemize}
    \item $\mathcal{T} = \{t_1, t_2, \dots, t_m\}$ defines a finite set of object types establishing a single-inheritance hierarchical structure ($t_i \sqsubseteq t_j$).
    \item $\mathcal{P} = \{p_1, p_2, \dots, p_n\}$ represents a set of first-order predicate signatures $p_k(x_1 - t_1, \dots, x_r - t_r)$, formalizing state relations.
    \item $\mathcal{O} = \{o_1, o_2, \dots, o_p\}$ contains object instances instantiated within the typing hierarchy ($o_i \in t_j$).
    \item $\mathcal{A} = \{a_1, a_2, \dots, a_q\}$ constitutes a set of parameterised STRIPS action schemas $a = \langle \text{params}(a), \text{pre}(a), \text{eff}^+(a), \text{eff}^-(a) \rangle$, specifying parameter bindings, precondition logical formulas, add-effects, and delete-effects.
    \item $I \subseteq \mathcal{P}(\mathcal{O})$ denotes the set of grounded first-order state literals holding true at initial state $S_0$.
    \item $G \subseteq \mathcal{P}(\mathcal{O})$ represents the goal state conjunction to be satisfied at plan termination state $S_n$.
\end{itemize}

Despite decades of mathematical refinement and the development of highly optimized heuristic search planners—such as Fast Downward \cite{helmert2006}—\textbf{authoring, formalizing, debugging, and maintaining PDDL domain models remains a severe human engineering bottleneck} \cite{georgievski2024}. Translating human domain requirements or informal natural language operational guidelines into syntactically valid and planner-executable PDDL requires specialized expertise in mathematical logic and formal modeling. Even subtle human errors—such as an undeclared predicate parameter, a mismatched variable binding, a missing precondition predicate, or an invalid type inheritance declaration—can cause immediate parser collapse, state-space explosion, or goal reachability failure.

In recent years, the rapid advancement of code-pretrained Large Language Models (LLMs) such as GPT-4, Claude 3.5 Sonnet, and Gemini 1.5/2.5 has demonstrated remarkable capabilities in software engineering and natural language formalization \cite{chen2021, oswald2024}. This has inspired substantial research into utilizing LLMs as automated translators to convert natural language domain descriptions into formal PDDL models \cite{valmeekam2023, kambhampati2024}. However, Large Language Models operate fundamentally as statistical pattern-matching engines over word token probability distributions. They lack intrinsic symbolic reasoning capabilities, formal world models, and deterministic syntax verification engines \cite{lin2023}. Consequently, direct one-shot LLM translation frequently emits syntactically malformed or logically flawed PDDL code blocks that fail when passed to symbolic classical planners.

\section{Gaps in the State of Art}
A rigorous examination of existing literature reveals three fundamental architectural gaps in current LLM-driven planning domain generation frameworks:

\subsection{Gap 1: Syntax Collapse in Unconstrained Text Generation}
Existing domain generation pipelines typically prompt LLMs to emit raw Lisp-formatted PDDL strings directly in an unconstrained manner. Because LLMs generate text token by token without strict grammatical constraints, they frequently suffer from \textbf{syntax collapse}. Common syntax collapse failure modes include:
\begin{enumerate}[label=(\alph*)]
    \item Unmatched closing parentheses in nested Lisp S-expressions.
    \item Keyword hallucinations, such as generating \texttt{:preconditions} instead of \texttt{:precondition} or \texttt{:effect-list} instead of \texttt{:effect}.
    \item Invalid parameter arities, where action preconditions reference predicate parameters not declared in the action's \texttt{:parameters} signature.
    \item Broken type hierarchy signatures, such as declaring multiple parent types in violation of standard STRIPS single-type inheritance rules.
\end{enumerate}
Direct text generation fails to separate natural language semantic extraction from formal target syntax generation.

\subsection{Gap 2: Absence of Closed-Loop Fault Localization \& Targeted AST Repair}
When an unconstrained LLM generates a PDDL domain that fails execution under a symbolic planner or validator (such as \texttt{VAL} or \texttt{Fast Downward}), current baseline frameworks rely on naive "unstructured retry" re-prompting \cite{more2024}. In an unstructured retry loop, the raw error output string from the compiler or planner is appended to the prompt, and the LLM is instructed to "try again and fix the domain." 

This open-loop approach suffers from two severe defects:
\begin{enumerate}[label=(\alph*)]
    \item Raw error traces (e.g., C++ stack dumps or ungrounded state search failures) are extremely noisy and lack explicit diagnostic fault localization.
    \item Without a diagnostic triage mechanism that isolates the precise Abstract Syntax Tree (AST) node responsible for the failure (such as action \texttt{load-truck}'s parameter list), the LLM attempts to re-generate the entire domain file from scratch. This leads to collateral corruption of previously valid action schemas and predicate signatures, causing LLMs to spin in unguided, non-convergent retry loops.
\end{enumerate}

\subsection{Gap 3: Neglect of the Continuous KEPS Domain Lifecycle}
Existing literature treats planning domain generation as a static, one-shot translation task. However, in practical engineering environments, planning domain models are dynamic artifacts that undergo continuous requirement evolution and maintenance \cite{georgievski2023}. Georgievski \cite{georgievski2023, georgievski2024} established that AI planning engineering requires a complete 4-phase \textbf{Knowledge Engineering for Planning and Scheduling (KEPS)} lifecycle:
\begin{enumerate}[label=Phase \arabic*:]
    \item \textbf{Acquisition \& Conceptualization:} Extracting domain entities, relations, and operational rules from requirements.
    \item \textbf{Formulation \& Encoding:} Translating conceptual models into formal executable representations.
    \item \textbf{Verification \& Validation (V\&V):} Checking static syntax validity, predicate consistency, and dynamic plan solvability.
    \item \textbf{Maintenance \& Evolution:} Modifying domain models to accommodate new actions or predicates while guaranteeing backward compatibility with historical regression test suites.
\end{enumerate}
Current LLM tools focus exclusively on initial generation, completely ignoring Phase 4 Maintenance and regression testing during domain evolution.

\section{Research Questions}
To address these fundamental challenges and bridge the identified gaps in the state of the art, this thesis investigates the primary research question:
\begin{quote}
    \textbf{RQ1:} How can a closed-loop neuro-symbolic agent architecture automate the complete 4-phase KEPS lifecycle for PDDL domain modelling from natural language descriptions without requiring human-in-the-loop intervention?
\end{quote}

To systematically evaluate and answer \textbf{RQ1}, this research decomposes the primary investigation into three specific, measurable sub-questions:
\begin{itemize}
    \item \textbf{RQ1.1 (Intermediate AST Schemas vs. Unconstrained Text Generation):} How effective is enforcing an intermediate Pydantic Abstract Syntax Tree (AST) schema contract (\texttt{DomainBlueprint}) paired with a deterministic programmatic Lisp compiler in eliminating syntactic formatting errors compared to unconstrained direct LLM text generation?
    \item \textbf{RQ1.2 (Symbolic Fault Localization \& Targeted AST Repair):} How accurately can a diagnostic triage engine—combining symbolic binary Oracles (\texttt{VAL} validator and \texttt{Fast Downward} planner)—calculate Fault Suspiciousness Rank Scores ($0.0$ to $1.0$) to guide targeted, node-level AST repair loops?
    \item \textbf{RQ1.3 (Phase 4 Maintenance \& Historical Regression Preservation):} How does incorporating an automated Phase 4 historical regression suite preserve backward compatibility and prevent semantic regression during continuous domain evolution?
\end{itemize}

\section{Key Technical Contributions}
This thesis resolves these research questions through the design, implementation, and empirical validation of a **Unified Autonomous Lifecycle Agent for PDDL Domain Modelling**. The key technical contributions are:
\begin{enumerate}
    \item \textbf{A Unified 4-Phase KEPS Lifecycle Agent Architecture:} Implemented an autonomous neuro-symbolic state-machine agent using \textbf{LangGraph}, establishing centralized state management (\texttt{AgentState}), persistent execution memory, and deterministic conditional routing edges across Acquisition, Encoding, V\&V, and Maintenance.
    \item \textbf{Intermediate Relational AST Schemas (\texttt{DomainBlueprint}, \texttt{ProblemBlueprint}):} Formulated strict Pydantic JSON contracts that enforce structural typing, predicate signatures, and action preconditions/effects on LLM outputs, coupled with a **deterministic programmatic Lisp compiler** that programmatically guarantees 100\% syntactically sound PDDL code construction.
    \item \textbf{Dual Symbolic Binary Oracles \& Diagnostic Triage Engine:} Integrated \texttt{VAL} (syntax and plan validator) and \texttt{Fast Downward} (heuristic search planner) into an automated triage engine (\texttt{nodes/triage.py}) that parses execution traces, categorizes failure modes (\texttt{SYNTAX}, \texttt{TYPE\_SIGNATURE}, \texttt{GROUNDING\_UNSOLVABILITY}), and computes Fault Suspiciousness Rank Scores for surgical AST patching.
    \item \textbf{Phase 4 Historical Regression Maintenance Suite:} Implemented a continuous maintenance module (\texttt{nodes/maintenance.py}) that automatically validates modified domain models against historical problem instances to preserve backward compatibility during domain evolution.
    \item \textbf{\texttt{PDDLLifeCycleBench} Dataset \& Interactive Web Application:} Established a comprehensive benchmark suite of 15 tasks across 5 classical IPC domain families (\texttt{Blocksworld}, \texttt{Logistics}, \texttt{Depots}, \texttt{Satellite}, \texttt{Zenotravel}) and deployed an interactive **Streamlit** visual web application for real-time thought-streaming and state graph inspection.
\end{enumerate}

\section{Thesis Structure}
The remainder of this report is structured into five detailed chapters as follows:
\begin{itemize}
    \item \textbf{Chapter 2 (Background \& Related Work):} Provides a thorough literature review covering classical AI planning fundamentals, knowledge engineering lifecycles (KEPS), LLM formalization limitations, neuro-symbolic LLM-Modulo architectures, and automated program repair frameworks.
    \item \textbf{Chapter 3 (System Architecture \& The 4-Phase KEPS Lifecycle):} Details the overall system architecture, the LangGraph state machine design, and the step-by-step operational mechanics across all 4 KEPS lifecycle phases.
    \item \textbf{Chapter 4 (Technical Implementation):} Explains the Pydantic AST schema definitions, the programmatic Lisp compiler algorithms, the multi-provider LLM integration layer, symbolic binary execution wrappers, and the Streamlit visual dashboard.
    \item \textbf{Chapter 5 (Empirical Evaluation \& Results):} Presents quantitative empirical evaluations on the \texttt{PDDLLifeCycleBench} benchmark dataset, comparing our closed-loop AST agent against \texttt{ZERO\_SHOT} and \texttt{UNSTRUCTURED\_RETRY} baselines, alongside qualitative fault repair case studies.
    \item \textbf{Chapter 6 (Conclusion \& Future Directions):} Summarizes the primary technical contributions, synthesizes research findings, and outlines promising future directions for temporal, numeric, and multi-agent domain engineering.
\end{itemize}


% --- Chapter 2: Background & Related Work ---
\chapter{Background \& Related Work}

\section{Knowledge Engineering for Planning \& Scheduling (KEPS) \& The PDDL Lifecycle}
Automated planning is a central field of Artificial Intelligence dedicated to synthesizing sequences of actions (plans) that transition an agent from an initial world state $S_0$ into a goal state $S_g$ \cite{ghallab2004}. Formally, a classical planning task is represented as a tuple $\Pi = \langle \mathcal{T}, \mathcal{P}, \mathcal{O}, \mathcal{A}, I, G \rangle$, where $\mathcal{T}$ is a set of object types, $\mathcal{P}$ is a set of first-order predicate signatures, $\mathcal{O}$ is a set of object instances, $\mathcal{A}$ is a set of action schemas with preconditions and state transition effects, $I$ is the initial state, and $G$ is the goal predicate set.

The Planning Domain Definition Language (PDDL) standardizes this representation by decoupling world dynamics into \texttt{domain.pddl} and instance specifics into \texttt{problem.pddl} \cite{mcdermott1998}. However, manual domain modeling remains an engineering bottleneck \cite{georgievski2024}. To address this, Georgievski \cite{georgievski2023, georgievski2024} established the 4-phase \textbf{Knowledge Engineering for Planning and Scheduling (KEPS)} lifecycle:
\begin{enumerate}
    \item \textbf{Phase 1 (Acquisition \& Conceptualization):} Eliciting knowledge and mapping domain concepts.
    \item \textbf{Phase 2 (Formulation \& Encoding):} Synthesizing formal domain representations.
    \item \textbf{Phase 3 (Verification \& Validation - V\&V):} Checking static syntax and dynamic plan solvability.
    \item \textbf{Phase 4 (Maintenance \& Evolution):} Managing continuous updates and historical regression suites.
\end{enumerate}

\section{LLMs as Domain Model Generators \& Planning Formalizers}
With the rise of code-trained Large Language Models (LLMs), researchers have investigated automated natural-language-to-PDDL translation \cite{oswald2024, valmeekam2023}. Early work by Oswald et al. \cite{oswald2024} evaluated direct one-shot LLM translation (the \texttt{ZERO\_SHOT} paradigm), revealing that unconstrained text generation suffers from \textbf{syntax collapse}—generating unmatched parentheses, invalid colon keywords, or mismatched predicate arities.

Lin et al. \cite{lin2023} and Valmeekam \& Kambhampati \cite{valmeekam2023} further demonstrated that raw LLMs lack formal internal world models, failing to generate sound plans (<30\% validity rate) without external symbolic guidance. Additional studies on domain scaling and type consistency confirm that direct text generation struggles with complex multi-object IPC benchmark domains.

\section{Neuro-Symbolic AI \& LLM-Modulo Architectures}
To bridge the gap between neural generation and symbolic precision, Kambhampati et al. \cite{kambhampati2024} proposed the \textbf{LLM-Modulo Paradigm}. Rather than relying on LLMs as self-contained reasoning engines, LLM-Modulo architectures position LLMs as hypothesis generators operating under external symbolic verifiers (modulo critics).

In parallel, Agarwal \& Sreepathy \cite{agarwal2024} introduced the \textbf{Translate-Infer-Compile (TIC)} framework, showing that enforcing intermediate structured representations (such as Abstract Syntax Tree schemas) prior to target code compilation eliminates string formatting errors. Hybrid neuro-symbolic systems leverage LLM natural language flexibility while guaranteeing formal logical soundness.

\section{Symbolic Validation, Fault Localization \& Automated Model Repair}
Verification and validation (V\&V) rely on symbolic binary Oracles:
\begin{itemize}
    \item \textbf{VAL (Plan Validator):} Howey, Long, and Fox \cite{howey2004} developed VAL to verify first-order state transitions $S_{i+1} = (S_i \setminus \text{eff}^-(a_i)) \cup \text{eff}^+(a_i)$, checking variable bindings and precondition satisfaction.
    \item \textbf{Fast Downward Planner:} Helmert \cite{helmert2006} introduced Fast Downward, a domain-independent planner using heuristic state-space search ($A^*$, Fast-Forward) to verify goal reachability.
\end{itemize}

When verification fails, naive systems employ unstructured retry re-prompting \cite{more2024}. In contrast, automated program repair (APR) and fault localization literature \cite{autostructor2024, orcaloca2024} demonstrates that parsing verifier log traces into diagnostic fault suspiciousness scores enables surgical, component-level AST repair, preventing collateral code damage.

\section{Autonomous Software Engineering Agents \& Domain Maintenance/Evolution}
Recent surveys on autonomous agents \cite{wang2024} and software engineering agent state machines \cite{chen2021} justify state-graph architectures (such as LangGraph) with persistent memory for complex multi-step workflows.

Furthermore, work on domain model evolution \cite{plan2evolve2024} demonstrates that maintaining domains over time (\textbf{Phase 4 Maintenance}) requires continuous validation against historical problem instances to prevent regression breakages. Our thesis unplugs these independent threads into a single, unified 4-phase KEPS agent architecture.


% --- Bibliography ---
\begin{thebibliography}{99}
\bibitem{agarwal2024} Agarwal, K., \& Sreepathy, V. (2024). \textit{Translate-Infer-Compile (TIC) for accurate text-to-plan using LLMs and Logical Representations}. arXiv preprint arXiv:2403.02309.
\bibitem{autostructor2024} AutoStructor Team (2024). \textit{AutoStructor: A Generative AI-Based Framework for Automated Program Repair with Deep Learning-Guided Fault Localization}. arXiv preprint.
\bibitem{chen2021} Chen, M., et al. (2021). \textit{Evaluating Large Language Models Trained on Code}. arXiv preprint arXiv:2107.03374.
\bibitem{georgievski2023} Georgievski, I. (2023). \textit{Software Development Life Cycle for Engineering AI Planning Systems}. In Proceedings of ICAPS/KEPS, pp. 751–760.
\bibitem{georgievski2024} Georgievski, I. (2024). \textit{Engineering AI Planning Systems}. Habilitation Thesis, University of Stuttgart.
\bibitem{ghallab2004} Ghallab, M., Nau, D., \& Traverso, P. (2004). \textit{Automated Planning: Theory \& Practice}. Morgan Kaufmann Publishers.
\bibitem{helmert2006} Helmert, M. (2006). \textit{The Fast Downward Planning System}. Journal of Artificial Intelligence Research, 26, 191–246.
\bibitem{howey2004} Howey, R., Long, D., \& Fox, M. (2004). \textit{VAL: Automatic plan validation, continuous effects and mixed initiative planning using PDDL}. In Proceedings of IEEE ICTAI, pp. 294–301.
\bibitem{kambhampati2024} Kambhampati, S., Valmeekam, K., Guan, L., et al. (2024). \textit{LLMs Can’t Plan, But Can Help Planning in LLM-Modulo Frameworks}. arXiv preprint arXiv:2402.01817.
\bibitem{lin2023} Lin, S., et al. (2023). \textit{On the Planning Abilities of Large Language Models: A Comparative Study}. In Proceedings of AAAI 2023.
\bibitem{mcdermott1998} McDermott, D., et al. (1998). \textit{PDDL - The Planning Domain Definition Language}. Technical Report CVC TR-98-003, Yale University.
\bibitem{more2024} MoRe Team (2024). \textit{MoRe: LLM-Based Domain Model Generation with Hybrid Self-Refinement}. Pre-print.
\bibitem{orcaloca2024} OrcaLoca Team (2024). \textit{OrcaLoca: An LLM Agent Framework for Software Issue Localization}. arXiv preprint.
\bibitem{oswald2024} Oswald, J., Srinivas, K., Kokel, H., Lee, J., Katz, M., \& Sohrabi, S. (2024). \textit{Large Language Models as Planning Domain Generators}. In Proceedings of ICAPS 2024. arXiv preprint arXiv:2405.06650.
\bibitem{plan2evolve2024} plan2evolve Team (2024). \textit{plan2evolve: LLM Self-Evolution for Improved Planning Capability via Automated Domain Generation}. arXiv preprint.
\bibitem{valmeekam2023} Valmeekam, K., \& Kambhampati, S. (2023). \textit{Large Language Models as Zero-Shot Planners: Revisiting Plans, Goals, and Hyperactions}. arXiv preprint arXiv:2305.09626.
\bibitem{wang2024} Wang, L., et al. (2024). \textit{A Survey on Large Language Model-based Autonomous Agents}. IEEE Transactions on Knowledge and Data Engineering.
\end{thebibliography}

\end{document}
